\section{Unified Document Parsing and Citation Grounding}
\label{sec:parsing}

Scientific full texts arrive in heterogeneous formats: structured JATS~XML from PubMed Central or scholarly PDFs from OpenAlex and manual curation. To ensure downstream extraction prompts and grounding evaluators operate identically across formats, \texttt{parsers.py} and \texttt{table\_renderer.py} convert all documents into a unified, citation-grounded \texttt{Node} object tree.

\subsection{Format Auto-Detection and GROBID Conversion}
\texttt{parse\_document(content, format="auto", ...)} automatically dispatches input payloads:
\begin{enumerate}
    \item \textbf{PDF Input:} Detected via \texttt{\%PDF} magic bytes or \texttt{.pdf} suffix. \texttt{GrobidClient} submits the PDF to a GROBID~\cite{grobid} REST service (\texttt{/api/processFulltextDocument}) to produce TEI~XML. When table or figure cropping is enabled, \texttt{tei\_coordinates=["figure"]} is requested so every table and figure carries PDF page bounding boxes (\texttt{coords}).
    \item \textbf{TEI~XML Input:} Detected via \texttt{<TEI}, \texttt{<teiHeader}, or the TEI namespace.
    \item \textbf{JATS~XML Input:} Default XML fallback for PMC full-text articles.
\end{enumerate}

\subsection{Hierarchical \texttt{Node} Tree and Sentence Segmentation}
\texttt{parse\_xml()} traverses the \texttt{<body>} (or \texttt{<text>}) element using \texttt{BeautifulSoup(..., "xml")}, naturally excluding reference lists (\texttt{<ref-list>}), acknowledgments, and boilerplate to preserve LLM context budget. Every element receives a deterministic hierarchical ID via \texttt{construct\_node\_id()}:
\begin{itemize}
    \item \texttt{RootNode}: Root container for \texttt{<body>} / \texttt{<text>}.
    \item \texttt{SectionNode} (\texttt{sec-\{i\}}): Matches \texttt{<sec>} and \texttt{<div>} elements with section headings.
    \item \texttt{ParagraphNode} (\texttt{p-\{j\}}): Matches \texttt{<p>} elements (with nested tables stripped to avoid duplication).
    \item \texttt{SentenceNode} (\texttt{s-\{k\}}): Generated by \texttt{add\_sentences()} using the biomedical spaCy model \texttt{en\_core\_sci\_sm}~\cite{scispacy2019} (static vectors disabled for speed).
    \item \texttt{TableNode} (\texttt{table-\{m\}}): Matches \texttt{<table-wrap>}, \texttt{<table>}, and \texttt{<figure type="table">}, storing \texttt{caption}, \texttt{columns}, \texttt{rows}, \texttt{footnotes}, grid spans (\texttt{colspan}/\texttt{rowspan}), \texttt{coords}, and optional rendered PNG \texttt{image\_bytes}.
    \item \texttt{ImageNode} (\texttt{fig-\{n\}}): Matches non-table \texttt{<fig>} / \texttt{<figure>} elements with captions, \texttt{coords}, and optional PNG \texttt{image\_bytes}.
\end{itemize}

This structure produces unambiguous citation tags such as \texttt{[sec-0/p-1/s-2]} or \texttt{[sec-2/sec-0/table-0/row-3]}. During evaluation, \texttt{get\_node\_by\_id(tree, tag)} resolves any cited tag back to its exact \texttt{Node} in $O(\text{depth})$ time.

\subsection{PDF Table and Figure Cropping}
Text extraction from complex multi-column PDF tables frequently loses cell alignment or merges footnotes. Using PyMuPDF (\texttt{fitz}) and Pillow (\texttt{table\_renderer.py}), \texttt{crop\_tables\_from\_pdf()} and \texttt{crop\_figures\_from\_pdf()} parse GROBID's \texttt{coords} strings (\texttt{page,x,y,w,h;...}) and render high-resolution PNG crops at $200\;\text{DPI}$ with $8.0\;\text{pt}$ padding:
\begin{itemize}
    \item \textbf{Tables:} Cropped across the full page width from \texttt{table\_y0 - padding} down to the next table's top or page bottom, preventing clipping of wide columns and table footnotes. Multi-page tables are vertically stitched into a single PNG.
    \item \textbf{Figures:} Cropped tightly around the GROBID figure bounding box, with multi-region stitching if needed.
\end{itemize}

\subsection{Serialization for LLMs}
The sentence-segmented tree is serialized into two formats:
\begin{itemize}
    \item \textbf{Text-Only (\texttt{to\_llm\_string}):} Produces an indented, bracket-tagged plain-text representation suitable for text-only LLMs and JATS~XML inputs.
    \item \textbf{Multimodal Content Blocks (\texttt{to\_llm\_blocks}):} Produces LangChain-compatible \texttt{HumanMessage} content blocks. When \texttt{image\_bytes} are present on a \texttt{TableNode} or \texttt{ImageNode}, the serializer inserts a location anchor (\texttt{>>> [TABLE IMAGE ATTACHED: sec-1/table-0] <<<}) followed by a base64 PNG \texttt{image\_url} block---either \texttt{interleaved} inline or grouped in an \texttt{appendix} section. If \texttt{image\_bytes} is absent (e.g., JATS~XML), it automatically falls back to pipe-separated text tables.
\end{itemize}
